\section{Conclusion}\label{sec:conclusion}
In the matched NLP-assisted GDP implementation studied here, radial separation gives a small but repeatable computational benefit over point tangents. It generates fewer cuts on average but spends substantially more recorded row-separation time per run. The advantage depends on the formulation, tree policy, shared initialization and tested instances, and it does not carry over to higher external accepted counts. Supported exact conic alternatives have lower common-solved mean times in these comparisons, while external coverage is mixed.

The mathematical analysis states what holds independently of these timings: valid perspective inequalities, finite fixed-residual separation under explicit compactness and enforcement conditions, and a safe residual-calibrated treatment of small weights. The diagnostic and counterexamples show that these guarantees imply neither universal cut dominance nor a general linear objective-error bound without additional joint regularity. With the complete model specifications, retained negative outcomes, original-model validation and reproducible artifacts, the study gives a concrete basis for deciding whether radial separation is worth its cost in a particular GDP solver.
